% =============================================================================
%  Notas de clase — Sesión 12: Torres de Hanoi en Arte ASCII
%  Programación Avanzada — Universidad Panamericana
%  Compilar: pdflatex sesion12_hanoi_ascii.tex
% =============================================================================
\documentclass[12pt, oneside]{article}
\usepackage{progavanzada}

\begin{document}

\portada{Sesión 12}{Torres de Hanoi en Arte ASCII}{2026}

\tableofcontents
\newpage

% =============================================================================
\section{La leyenda y el problema}
% =============================================================================

En 1883, el matemático francés \textbf{Édouard Lucas} lanzó al mercado un
juguete llamado \textit{La Tour d'Hanoï} — un rompecabezas de ocho discos
de madera de distinto tamaño apilados en una de tres estacas.
Para justificar el precio, lo acompañó con una leyenda apócrifa:

\begin{quote}
\textit{En un templo de Benarés, en la India, los sacerdotes del dios Brahma
mueven en silencio 64 discos de oro puro entre tres agujas de diamante,
una y otra vez, sin descanso.  Cuando terminen de moverlos todos a la
tercera aguja, el mundo llegará a su fin.}
\end{quote}

Lucas la inventó él mismo para darle misterio al juguete, pero la pregunta
matemática que esconde es completamente real:
\textbf{¿cuántos movimientos mínimos se necesitan para mover $n$ discos?}

\begin{tecnico}[¿Cuándo termina el mundo?]
  Para $n = 64$ discos, el mínimo de movimientos es $2^{64} - 1
  \approx 1.8 \times 10^{19}$.
  Si los sacerdotes movieran un disco por segundo, sin parar, tardarían
  aproximadamente \textbf{585 mil millones de años} — unas 42 veces la edad
  actual del universo.  El mundo, al parecer, tiene tiempo.\\
  \textit{Referencia:} Lucas, É. \textit{Récréations Mathématiques}.
  Gauthier-Villars, 1883.
\end{tecnico}

% =============================================================================
\section{Las reglas del juego}
% =============================================================================

El problema tiene exactamente \textbf{tres reglas}:

\begin{enumerate}
  \item Solo puedes mover \textbf{un disco a la vez}.
  \item Solo puedes mover el disco que está en el \textbf{tope} de una torre.
  \item \textbf{Nunca} puedes colocar un disco sobre uno más pequeño.
\end{enumerate}

El objetivo es trasladar todos los discos de la torre \textbf{A} a la torre
\textbf{C}, usando la torre \textbf{B} como auxiliar, respetando las tres
reglas en todo momento.

\begin{verbatim}
Estado inicial:        Estado final:

    |      |      |        |      |      |
   ===     |      |        |      |     ===
  =====    |      |        |      |    =====
 =======   |      |        |      |   =======
=========  |      |        |      |  =========
    A      B      C            A      B      C
\end{verbatim}

% =============================================================================
\section{¿Cuántos movimientos se necesitan?}
% =============================================================================

Antes de programar nada, vale la pena entender cuánto trabajo implica el
problema.  Para $n$ discos, el \textbf{mínimo de movimientos} es:

\[
  M(n) = 2^n - 1
\]

\subsection{¿De dónde viene esa fórmula?}

Piensa recursivamente:

\begin{itemize}
  \item Para mover el disco más grande (el $n$) de A a C, necesitas
        que los $n-1$ discos que están encima \textbf{ya estén en B}.
  \item Mover esos $n-1$ discos de A a B cuesta $M(n-1)$ movimientos.
  \item Mover el disco $n$ de A a C cuesta \textbf{1} movimiento.
  \item Mover los $n-1$ discos de B a C cuesta $M(n-1)$ movimientos.
\end{itemize}

La recurrencia es:
\[
  M(1) = 1 \qquad M(n) = 2\,M(n-1) + 1
\]

Resolviendo:
\[
  M(n) = 2^n - 1
\]

\begin{center}
\begin{tabular}{ccr}
  \toprule
  \textbf{Discos} & \textbf{$2^n - 1$} & \textbf{Movimientos} \\
  \midrule
  1 & $2^1 - 1$ & 1 \\
  2 & $2^2 - 1$ & 3 \\
  3 & $2^3 - 1$ & 7 \\
  4 & $2^4 - 1$ & 15 \\
  8 & $2^8 - 1$ & 255 \\
  64 & $2^{64} - 1$ & $\approx 1.8 \times 10^{19}$ \\
  \bottomrule
\end{tabular}
\end{center}

\begin{recuerda}
  La recurrencia $M(n) = 2\,M(n-1) + 1$ \textbf{es} la descripción del
  algoritmo recursivo que implementarás en la sesión 13.
  Hoy no lo programas: tu equipo mueve los discos a mano (el usuario
  escribe los movimientos).  En la sesión 13 harás que la computadora
  los calcule sola.
\end{recuerda}

\begin{tecnico}[La conexión con el código binario]
  Hay una propiedad fascinante: en la secuencia óptima de movimientos,
  el disco que se mueve en el paso $k$ es el mismo que el bit menos
  significativo diferente de cero en la representación binaria de $k$.

  Para $n=3$, los primeros movimientos son:
  \begin{center}
  \begin{tabular}{ccc}
    \toprule
    \textbf{Paso $k$} & \textbf{$k$ en binario} & \textbf{Disco} \\
    \midrule
    1 & \texttt{001} & disco 1 (bit 0) \\
    2 & \texttt{010} & disco 2 (bit 1) \\
    3 & \texttt{011} & disco 1 (bit 0) \\
    4 & \texttt{100} & disco 3 (bit 2) \\
    5 & \texttt{101} & disco 1 (bit 0) \\
    6 & \texttt{110} & disco 2 (bit 1) \\
    7 & \texttt{111} & disco 1 (bit 0) \\
    \bottomrule
  \end{tabular}
  \end{center}
  Esta propiedad conecta el problema con el \textit{código Gray}
  (código binario reflejado), usado en electrónica para evitar glitches
  al incrementar contadores.  El espacio de estados de Hanoi también
  tiene la forma del \textbf{triángulo de Sierpiński} — uno de los
  fractales más famosos de las matemáticas.\\
  \textit{Referencia:} Gardner, M. \textit{Mathematical Games}.
  Scientific American, 1958.
\end{tecnico}

% =============================================================================
\section{Arte ASCII como interfaz}
% =============================================================================

El programa que van a construir hoy no usa ventanas ni botones:
toda su interfaz vive en la terminal, hecha de caracteres.
Eso es \textbf{arte ASCII aplicado} a la programación.

\subsection{¿Por qué importa el diseño visual del programa?}

Un programa que funciona pero es difícil de leer es un programa incompleto.
La claridad visual de la salida es parte de la calidad del código.
Cuando la interfaz de una herramienta de terminal es buena, el usuario
sabe inmediatamente qué está pasando sin leer manual.

Aquí está la diferencia entre un estado impreso de forma mínima y
uno bien diseñado:

\begin{verbatim}
Sin diseño:              Con diseño:
A: [4,3]                     |          |          |
B: [2]                      ===         |          |
C: [1]                    =====         |          |
                          =======      ===         |
                          =========   =====       ===
                          =========  =========  =========
                               A          B          C
\end{verbatim}

Ambos contienen la misma información, pero el segundo la hace
inmediatamente comprensible.

\subsection{Algunos caracteres útiles para diseño en terminal}

\begin{center}
\begin{tabular}{llp{6cm}}
  \toprule
  \textbf{Categoría} & \textbf{Caracteres} & \textbf{Uso típico} \\
  \midrule
  Líneas simples    & \texttt{| - =}              & Postes y bases básicas \\
  Bloques de sombra & \texttt{U+2591..2588} & Discos rellenos con gradiente \\
  Líneas dobles     & \texttt{U+2550..256C} & Marcos y bordes elegantes \\
  Flechas           & \texttt{U+2192..21D0} (\texttt{->}, \texttt{<-}) & Indicadores de movimiento \\
  Figuras geométricas & \texttt{U+25B2,25BC..} (\texttt{\^{} v *}) & Decoración de discos \\
  Relleno alternativo & \texttt{\# * \textasciitilde{} \^{}}    & Diseños personalizados \\
  \bottomrule
\end{tabular}
\end{center}

\begin{consejo}
  Los caracteres Unicode de bloques (\texttt{U+2591..2588},
  \texttt{U+2550..256C}, etc.) funcionan en la mayoría de terminales
  modernas (macOS, Linux y Windows Terminal).
  Si tu terminal no los muestra bien, usa los caracteres ASCII básicos
  (\texttt{= | \#}).
  Antes de diseñar, comprueba que tu terminal soporte lo que quieres usar.
\end{consejo}

% =============================================================================
\section{La actividad de hoy}
% =============================================================================

Trabajan en \textbf{equipos de tres personas}.
El objetivo es construir un programa que muestre las Torres de Hanoi
en arte ASCII y permita al usuario jugar.

\subsection{Lo que deben implementar}

\begin{enumerate}
  \item \textbf{Función de dibujo.}
        Dado el estado de las tres torres (qué discos tiene cada una),
        imprimir la representación visual en la consola.

  \item \textbf{Función de movimiento.}
        Recibir una torre de origen y una de destino, validar que el
        movimiento sea legal y ejecutarlo.

  \item \textbf{Interfaz de consola.}
        Un ciclo donde el usuario escribe el origen y destino de cada
        movimiento, el programa dibuja el estado actualizado y detecta
        cuando el jugador ganó.
\end{enumerate}

\subsection{Lo que deben entregar en Canvas}

\begin{itemize}
  \item El \textbf{código fuente} (archivo \texttt{.cpp}) con los
        tres integrantes del equipo en un comentario al inicio.
  \item \textbf{Capturas de pantalla o fotos} mostrando las torres en
        al menos \textbf{4 estados distintos}:
        \begin{enumerate}
          \item Estado inicial (todos los discos en la torre A).
          \item Un estado intermedio con los discos distribuidos.
          \item Un estado casi final.
          \item Estado final (todos los discos en la torre C).
        \end{enumerate}
\end{itemize}

\begin{advertencia}
  \textbf{No implementen el algoritmo que resuelve Hanoi automáticamente.}
  Eso es el tema de la sesión 13.  Hoy, los movimientos los elige el
  usuario escribiendo origen y destino.  Si el programa se resuelve solo,
  no cumple con el objetivo de la actividad.
\end{advertencia}

\subsection{Criterios de evaluación}

\begin{center}
\begin{tabular}{lp{7cm}c}
  \toprule
  \textbf{Criterio} & \textbf{Descripción} & \textbf{Peso} \\
  \midrule
  Funcionalidad      & El programa dibuja las torres correctamente y
                       las actualiza tras cada movimiento & 40\% \\
  Validación         & Rechaza movimientos ilegales con un mensaje claro & 20\% \\
  Condición de victoria & Detecta cuando todos los discos están en C & 10\% \\
  Diseño visual      & Las torres son claras, el arte es consistente
                       y estéticamente agradable & 20\% \\
  Legibilidad del código & Nombres de funciones y variables descriptivos,
                           estructura clara & 10\% \\
  \bottomrule
  & \textbf{Total} & 100\% \\
  \bottomrule
\end{tabular}
\end{center}

\begin{consejo}
  El \textbf{código ganador del concurso de la sesión 13} se usará como
  base para implementar el algoritmo recursivo.  Eso significa que vale
  la pena que el código esté bien estructurado desde hoy — no solo bonito
  visualmente.  Un diseño elegante con código ininteligible no será fácil
  de extender en la siguiente sesión.
\end{consejo}

% =============================================================================
\section{Análisis previo: ejercicios a mano}
% =============================================================================

Antes de ponerte a programar, resuelve estos ejercicios con lápiz y papel.
Te ayudarán a entender el problema desde adentro.

\begin{enumerate}
  \item Dibuja en papel el estado de las torres después de cada uno de los
        7 movimientos óptimos para mover 3 discos de A a C.
        Numera cada estado del 0 (inicial) al 7 (final).

  \item Verifica con la fórmula que $M(3) = 7$.
        Luego aplica la recurrencia para calcular $M(4)$ y $M(5)$.

  \item Indica cuántos frames se apilarían en el call stack cuando el
        algoritmo recursivo calcule la solución para $n = 4$.
        \textit{(Pista: la profundidad del árbol de recursión es $n$.)}

  \item ¿Cuál es el disco que siempre se mueve en el paso $k$
        si $k$ es potencia de 2?  Justifica con la tabla de la sección 3.

  \item \textbf{Desafío:}
        Con 4 agujas en lugar de 3, ¿necesitarías más o menos movimientos
        para mover $n$ discos?  Razona sin calcular.
\end{enumerate}

% =============================================================================
\section{Resumen}
% =============================================================================

\begin{recuerda}
\begin{itemize}
  \item Las Torres de Hanoi fueron inventadas por Édouard Lucas en 1883
        como juguete comercial; la leyenda del templo es apócrifa.
  \item Tres reglas: un disco a la vez, solo el tope, nunca un disco
        grande sobre uno pequeño.
  \item Mínimo de movimientos: $M(n) = 2^n - 1$.
        Recurrencia: $M(n) = 2\,M(n-1) + 1$.
  \item La secuencia óptima está relacionada con el código binario y
        el triángulo de Sierpiński.
  \item Hoy implementas la interfaz visual y los movimientos del usuario.
        La sesión 13 agrega el algoritmo recursivo que lo resuelve solo.
\end{itemize}
\end{recuerda}

\end{document}
